\section{CUDA 的主机-设备执行模型}

\subsection{一个源程序，两类执行位置}

CUDA 程序是\term{异构应用}{heterogeneous application}：同一个工程同时包含在 CPU 上运行的\term{主机代码}{host code}，以及在 GPU 上运行的\term{设备代码}{device code}。主机代码通常负责串行或并行度较低的工作，例如输入输出、内存管理与核函数调度；设备代码负责数据量大、并行度高的计算阶段。

\slidefig{0.94}{page-07.png}{CUDA/OpenCL 的集成式主机-设备执行模型}

一个典型的时间顺序可以写成：
\[
\text{主机准备数据}
\rightarrow
\text{启动核函数 A}
\rightarrow
\text{主机继续执行或等待}
\rightarrow
\text{启动核函数 B}
\rightarrow
\text{取回结果}.
\]

\subsection{核函数与单程序多数据模型}

\term{核函数}{kernel}是由许多 GPU 线程执行的设备函数。一次核函数启动（kernel launch）创建一个线程网格。网格中的所有线程运行同一份核函数代码，但每个线程使用自己的索引选择不同数据。这种编程方式称为\term{单程序多数据}{Single Program Multiple Data, SPMD}。

SPMD 的“同一程序”不意味着所有线程在每一时刻必然完成相同工作。线程可以依据索引进行条件判断；只是它们的代码文本相同。最简单的向量加法核函数逻辑是
\[
C[i]=A[i]+B[i],
\]
其中线程索引 $i$ 决定该线程访问哪一对输入元素。

\begin{examplebox}[title={主机与设备的职责边界}]
对向量加法而言，主机代码负责准备 $A$、$B$、$C$，把输入移动到设备可访问的位置，配置线程组织并启动核函数；设备线程只负责执行元素级加法。把数据移动、调度和计算区分开，是阅读任何 CUDA 程序的第一步。
\end{examplebox}
